% Exported from Math++. Compile with pdflatex.
\documentclass[11pt,a4paper]{article}
\usepackage[T1]{fontenc}
\usepackage[utf8]{inputenc}
\usepackage{lmodern}
\usepackage{amsmath,amssymb,amsthm,mathtools}
\usepackage{graphicx,booktabs,xcolor}
\usepackage[hidelinks]{hyperref}
\usepackage[margin=25mm]{geometry}
\theoremstyle{plain}
\newtheorem{theorem}{Theorem}
\newtheorem{lemma}[theorem]{Lemma}
\newtheorem{corollary}[theorem]{Corollary}
\newtheorem{proposition}[theorem]{Proposition}
\theoremstyle{definition}
\newtheorem{definition}[theorem]{Definition}
\newtheorem{example}[theorem]{Example}
\theoremstyle{remark}
\newtheorem*{remark}{Remark}
\setlength{\parindent}{0pt}\setlength{\parskip}{0.6em}
\title{Graphs in three dimensions}
\date{}
\begin{document}
\maketitle
Give a function of two variables and a range for each, and you get a surface. Drag to turn it,
scroll to zoom. The surface is coloured by height and shaded by a light.

\section*{Surfaces}

\begin{figure}[htbp]
\centering
\includegraphics[width=0.8\linewidth]{graphs-3d-fig1.png}
\caption{A saddle.}
\end{figure}
\begin{figure}[htbp]
\centering
\includegraphics[width=0.8\linewidth]{graphs-3d-fig2.png}
\caption{Ripples spreading from the origin.}
\end{figure}
\begin{equation*}
  f(x, y) \coloneqq  x e^{-x^2 - y^2}
\end{equation*}
\begin{figure}[htbp]
\centering
\includegraphics[width=0.8\linewidth]{graphs-3d-fig3.png}
\caption{A function you defined.}
\end{figure}
\section*{A slider turns the surface}

\begin{equation*}
  k \coloneqq  1, \; k \in [0.5, 3]
\end{equation*}
\begin{figure}[htbp]
\centering
\includegraphics[width=0.8\linewidth]{graphs-3d-fig4.png}
\caption{Move $k$ to change the number of waves.}
\end{figure}
\section*{Curves in space}

\begin{figure}[htbp]
\centering
\includegraphics[width=0.8\linewidth]{graphs-3d-fig5.png}
\caption{A helix.}
\end{figure}
\section*{Parametric surfaces}

\begin{figure}[htbp]
\centering
\includegraphics[width=0.8\linewidth]{graphs-3d-fig6.png}
\caption{A sphere.}
\end{figure}
\begin{figure}[htbp]
\centering
\includegraphics[width=0.8\linewidth]{graphs-3d-fig7.png}
\caption{A torus.}
\end{figure}
\end{document}
